\\section{结论}
\\label{sec:conclusion}

长期记忆有两个相互关联的问题：如何在实体和事实不断变化时保存不断增长的记录，以及如何在上下文窗口有限时只读取当前问题需要的证据。Deep-Dream通过保留文档片段作为事实记录，并把概念、关系、观测、版本和候选身份作为可修订的访问结构，回答了记忆如何存储的问题。读取时，智能体沿着这些结构打开相应的原文，比较早期和后续观测，并在证据足够时停止。

这种分工让每一部分承担清晰的责任。图结构缩小搜索范围，但本身不授权答案；原文片段提供具体措辞，但单独不能解决身份判断；智能体决定是否还需要查看另一份文档、关系或版本。MAB中的结果因此体现了完整读取路径的差异，而不是单独的图结构或原文消融。写入实验也表明，延迟对齐和并行处理可以降低构建成本并减少重复实体族，但目前没有测得一致的对齐质量提升。全文上下文在匹配题目上仍然更强，而MEME结果说明，保留历史记录并不等于已经规定了何时应当抑制历史。

正文中的实际轨迹进一步说明了这种分工：概念和关系负责定位候选，原文片段提供计算所需的具体数值。

本文的核心贡献，是把记忆存储和记忆读取分成两个可以分别改进、分别审计的问题：原文保留事实，概念、关系和版本提供可修订的索引，智能体自主决定读取路径和深度。后续研究可以在此基础上标注身份判断，使用相同资源评估蕴含和停止策略，并加入明确的当前状态、历史状态和删除抑制规则。
